\\
مسألهٔ مقدار اولیهٔ زیر را در نظر بگیرید:
\[
 \begin{cases}
 y'=x^2+y^2,\\
 y(0)=1.
 \end{cases}
 \qquad h=0.1
\]
\begin{enumerate}
 \item مقدار تقریبی \(y(0.2)\) را با روش اویلر محاسبه کنید.
 \item مقدار تقریبی \(y(0.2)\) را با روش رانگ--کوتای مرتبهٔ دوم محاسبه کنید.
 \item مقادیر آغازین لازم را با روش کلاسیک رانگ--کوتای مرتبهٔ سوم تولید کنید؛
 مقدار تقریبی \(y(0.2)\) تولیدشده را گزارش دهید و سپس \(y(0.3)\) را با روش سه‌گام آدامز--بشفورث محاسبه کنید.
\end{enumerate}

\begingroup\color{blue}
حل

تابع شیب برابر است با
\[
 f(x,y)=x^2+y^2.
\]

\begin{enumerate}
\item
در روش اویلر،
\begin{align*}
 y_1&=y_0+h f(x_0,y_0)=1+0.1(1)=1.1,\\
 y_2&=y_1+h f(x_1,y_1)
     =1.1+0.1\left(0.1^2+1.1^2\right)=1.222.
\end{align*}
بنابراین
\[
 \boxed{y(0.2)\approx1.222}.
\]

\item
طبق تعریف درس، منظور از رانگ--کوتای مرتبهٔ دوم در این سؤال روش هیون است:
\[
 K_1=f(x_n,y_n),\qquad
 K_2=f(x_n+h,y_n+hK_1),\qquad
 y_{n+1}=y_n+\frac{h}{2}(K_1+K_2).
\]
در گام نخست،
\[
 K_1=1,\qquad K_2=f(0.1,1.1)=1.22,
 \qquad y_1=1+\frac{0.1}{2}(1+1.22)=1.111.
\]
در گام دوم،
\begin{align*}
 K_1&=f(0.1,1.111)=1.244321,\\
 y_1+hK_1&=1.2354321,\\
 K_2&=f(0.2,1.2354321)=1.56629247371041,\\
 y_2&=1.111+\frac{0.1}{2}
       \left(1.244321+1.56629247371041\right)\\
    &=1.2515306736855205.
\end{align*}
پس
\[
 \boxed{y(0.2)\approx1.251530674}.
\]

\item
برای تولید مقادیر آغازین از روش کلاسیک رانگ--کوتای مرتبهٔ سوم استفاده می‌کنیم:
\begin{align*}
 k_1&=h f(x_n,y_n),\\
 k_2&=h f\left(x_n+\frac h2,y_n+\frac{k_1}{2}\right),\\
 k_3&=h f(x_n+h,y_n-k_1+2k_2),\\
 y_{n+1}&=y_n+\frac{k_1+4k_2+k_3}{6}.
\end{align*}
برای دو گام آغازین داریم
\[
\begin{array}{c|ccc|c}
 n&k_1&k_2&k_3&y_{n+1}\\ \hline
 0&0.1&0.1105&0.1266641&1.111444016666667\\
 1&0.124530780218413&0.140009377155417&0.164511666920879
   &1.252957342626827
\end{array}
\]
بنابراین
\[
 \boxed{y(0.2)\approx1.252957343}.
\]
اکنون
\begin{align*}
 f_0&=1,\\
 f_1&=0.1^2+(1.111444016666667)^2
      =1.245307802184134,\\
 f_2&=0.2^2+(1.252957342626827)^2
      =1.609902102442479.
\end{align*}
فرمول سه‌گام آدامز--بشفورث می‌دهد
\begin{align*}
 y_3
 &=y_2+\frac{h}{12}(23f_2-16f_1+5f_0)\\
 &=1.252957342626827+\frac{0.1}{12}
   \left(23(1.609902102442479)-16(1.245307802184134)+5\right)\\
 &=1.437147538637084.
\end{align*}
پس
\[
 \boxed{y(0.3)\approx1.437147539}.
\]
\end{enumerate}
\endgroup
